\section{How strong must selection be?}\label{sec:strength}

\noindent Combining the rate audit with the information model gives an order-of-magnitude check. If the brain increase were
entirely a correlated response to selection on cognitive capacity, the required sex-averaged selection differential on
capacity is the brain rate divided by $r_gh_ch_b$~\cite{lande1979}. With $r_g=0.24$~\cite{jansen2020}, $h_c^2=0.5$, and
$h_b^2=0.8$, the requirement is $\ReqBase$ standard deviations per generation (range $\ReqLo$ to $\ReqHi$ over the grid in
Table~\ref{tab:envelope}, which also covers the possibility that $r_g$ is overstated~\cite{border2022}). On the supply side,
a status--reproduction correlation of 0.19~\cite{vonrueden2016} implies a selection differential on men's standing of
$0.19\sqrt I$, where $I$ is the opportunity for selection among men. If standing correlates with capacity at $\rho_{ZC}$,
only men's standing matters, and capacity affects reproduction only through standing, the sex-averaged differential on
capacity is $\tfrac12\rho_{ZC}\,0.19\sqrt I$.

\begin{table}[t]
\caption{Plausibility envelope: required versus available selection on capacity (E7).}\label{tab:envelope}
\footnotesize
\begin{tabularx}{\columnwidth}{@{}Xccc@{}}
\toprule
\tabinput{generated/tab_envelope.tex}
\bottomrule
\end{tabularx}
\tablenote{$I$ and $\rho_{ZC}$ are assumptions, not estimates. Claim level: plausibility bound.}
\end{table}

Under that mediation assumption, a weak standing--capacity correlation of 0.05 supplies roughly \MultLoA{} to \MultLoB{} times
the required selection, and 0.3 supplies \MultHiA--\MultHiB{} times (Table~\ref{tab:envelope}). A floor gives the same answer by
another route: if one or two men in ten are excluded on an assessed score, the intensity is 0.2 to 0.35
(Corollary~\ref{cor:floor}), and with a validity of only $R=0.04$ the sex-averaged differential on capacity is still several
times $\ReqBase$.

\emph{What this does and does not show.} It shows that the requirement is small: no strong effect is needed anywhere in the
chain, so the absence of strong contemporary associations would not refute the mechanism. It does not show that the channel
operated. The product of two correlations estimates $\Cov(C,W)$ only if standing carries the whole association; in general
$\Cov(C,W)=\rho_{ZC}\Cov(Z,W)+\E[\Cov(C,W\mid Z)]$, and the second term is unmeasured and can have either sign. The pooled
status correlation also mixes several outcome measures~\cite{vonrueden2016}. The comparison is a bound on what would be
needed, not an estimate of what happened, and it cannot by itself show that the budget and not the demand set the pace.

\emph{What is known about the two unknowns.} The opportunity for selection among men is modest in humans compared with
other mammals and varies widely among populations~\cite{ross2023,brown2009}, so the grid's range for $I$ is deliberately
wide. For $\rho_{ZC}$ there is no estimate from a natural-fertility population, only brackets. In modern organizations
intelligence correlates 0.21 to 0.27 with leadership~\cite{judge2004}. Leaders in the ethnographic record are described by
expertise, good counsel and what one synthesis calls neural capital~\cite{garfield2019}. Impressions formed from short
behavioral episodes predict measured intelligence~\cite{borkenau2004}. Against these, Hadza hunting reputations track measured
skills only noisily~\cite{stibbardhawkes2018}, and strength predicts reputation better than any other measured
trait~\cite{apicella2014}. Values of $\rho_{ZC}$ between 0.05 and 0.3 are consistent with all of this, and so is a value near
zero in settings where standing rests on force. The single most valuable new measurement is therefore $\rho_{ZC}$: the
correlation between socially assessed standing and independently measured cognitive capacity in natural-fertility
populations, estimated separately for standing won by dominance and standing conferred as prestige.

\section{Other hominins}\label{sec:hominins}

\noindent The fossil and genetic record of other late hominins bears on the mechanism in three ways (Fig.~\ref{fig:fossils}).

\emph{Neanderthals show that the result did not require our kind of society.} Their brains were as large as those of
contemporaneous \Hs{} (Section~\ref{sec:parallel}), and the western Eurasian series leading to them, which begins with early members of the lineage
at Sima de los Huesos~\cite{arsuaga2014}, rose at \SerEur{} log units per Myr, about the pace of the African series. They used fire habitually~\cite{roebroeks2011}, built structures deep in
Bruniquel Cave 176,000 years ago~\cite{jaubert2016}, had auditory capacities matched to speech like
ours~\cite{condevalverde2021}, and suffered cranial trauma at rates similar to Upper Paleolithic \Hs, higher in males in both
groups~\cite{beier2018}. Yet their communities were small. The parents of an Altai woman were related at the level of
half-siblings~\cite{prufer2014}; at Chagyrskaya Cave, long runs of homozygosity and low Y-chromosome diversity indicate a
small community, of the order of twenty individuals, linked to others mainly by the movement of women~\cite{skov2022}. Early Upper Paleolithic \Hs{} at
Sunghir, by contrast, show no close inbreeding and belonged to a wider mating network~\cite{sikora2017}. Any account in which
large, dense networks are \emph{necessary} for large brains is therefore contradicted (objection O25). The present account
does not need them. In a small community everyone observes everyone for years, so validity can be high; what a small
community limits is the number of candidates, and hence the opportunity for selection, and it raises drift. The model
therefore predicts a slower and noisier response in such populations, not none. This is a historical conjecture and cannot be
tested with present data.

\emph{The two lineages differ in organization more than in size.} Neanderthal endocasts are elongated where those of
\Hs{} became globular~\cite{neubauer2018}, with relatively smaller cerebellar hemispheres~\cite{kochiyama2018}, and by one
argument more of the volume was committed to vision and to a heavier body~\cite{pearce2013}. In the terms of
Proposition~\ref{prop:tissue}, similar tissue $B$ with different organization $K$. Size is the wrong variable for comparing
the cognition of the two; it remains the right variable for the energetic argument, since tissue is what costs.

\emph{Small-brained species show the gates.} \textit{H.~naledi} and \textit{H.~floresiensis} persisted with
australopith-sized brains while large-brained lineages lived elsewhere~\cite{berger2015,dirks2017,kubo2013,sutikna2016}. The island
form is the clearer case for a budget reading: its brain size is consistent with dwarfing from a larger-bodied
ancestor~\cite{kubo2013}. A demand-side
mechanism that acts only where the budget allows and where tissue pays is compatible with such exceptions, but so is any
account that includes costs; the exceptions test the budget, not the demand. They also carry a warning. To the extent that
small-brained hominins made tools and coped with the same world, capacity is not simply volume, and the coupling $a$ should
not be assumed large.

A sharper test separates two things that covary in our own history. Cumulative culture is expected to track the number of
people who can learn from one another~\cite{powell2009,hill2014}; tissue, on the present account, should track the energy
budget and the validity of assessment. Neanderthals, with small networks and large brains, show that tissue did not wait for
large networks; whether their cultural accumulation was correspondingly slower is debated (prediction P16).
